% vim: textwidth=99

\documentclass[nobib, a4paper, twoside, justified]{tufte-book}

\usepackage[utf8]{inputenc}
\usepackage[UKenglish]{babel}
\usepackage{csquotes}
\usepackage[style=verbose, autocite=footnote, backend=biber]{biblatex}

\title{Displaying Heraldic Blazons}
\author{William Mathewson}
\publisher{University of Edinburgh}
\date{January 2018}

\addbibresource{bibliography.bib}

%%%%
%%
%% LOCALS are in `tufte-book-local.tex'

\usepackage{graphicx} % allow embedded images
\setkeys{Gin}{width=\linewidth,totalheight=\textheight,keepaspectratio}
\graphicspath{{graphics/}} % set of paths to search for images
\usepackage{amsmath}  % extended mathematics
\usepackage{booktabs} % book-quality tables
\usepackage{units}    % non-stacked fractions and better unit spacing
\usepackage{multicol} % multiple column layout facilities
\usepackage{fancyvrb} % extended verbatim environments
\usepackage{amsmath}
\usepackage{xspace}
\usepackage{makeidx}
\usepackage{mathtools}

\usepackage{float}
\usepackage{subfig}

\usepackage[final]{pdfpages}

%% Tell hyperref packag it may break long URLs however it feels
\usepackage{hyperref}
\def\UrlBreaks{\do\/\do-}

\newcommand{\dcim}{\emph{DCIM}\@\xspace}
\newcommand{\ipam}{\emph{IPAM}\@\xspace}
\newcommand{\dcimipam}{\emph{DCIM \& IPAM}\@\xspace}
\newcommand{\desired}{\emph{desired}\@\xspace}
\newcommand{\operational}{\emph{operational}\@\xspace}
\newcommand{\wikid}{\textsc{WikiD}\@\xspace}
\newcommand{\pywikid}{\textit{pywikid}\@\xspace}
\newcommand{\tardis}{\textit{The Tardis Project}\@\xspace}
\newcommand{\note}{\par\noindent{\fontsize{10}{12}\itshape Note}\vspace{0.2em}\par}
\newcommand{\code}[1]{\texttt{#1}}
\newcommand{\pythonthree}{Python$3$\@\xspace}

\newcommand{\defeq}{\stackrel{\mathclap{\normalfont\mbox{def}}}{=}}
\newcommand{\namedmap}[1]{\stackrel{\mathclap{\mbox{\small#1}}}{\longmapsto}}

%% R_{x->y} style
\newcommand{\relate}[3]{\text{\itshape#1 }_{#2\ \mapsto\ #3}}

%% x --R--> y style
%\newcommand{\relate}[3]{#2\namedmap{#1}#3}

%% xRy style
%\newcommand{\relate}[3]{_{#2}#1_{#3}}

\newcommand{\citetodo}{\sidenote{\todo{[Citation Needed]}}\@\xspace}
\newcommand{\reftodo}{\sidenote{\todo{[Reference Needed]}}\@\xspace}
\newcommand{\prooftodo}{\sidenote{\todo{[Substantiation Needed]}}\@\xspace}
\newcommand{\todo}[1]{{\noindent\textcolor{Red}{\textit{\quad#1}}\par}}

%%  Service description table
%
\newcommand{\servicedesc}[4]{%
  \vspace{\baselineskip}%
  \noindent%
  {%
    \footnotesize%
    \begin{tabular}{llll}%
      \textbf{Service Name}  & \textbf{Type}  & \textbf{Graph} & \textbf{Integrates with}\\%
      #1                     & #2             & #4             & #3  \\%
    \end{tabular}%
  }%
  \vspace{\baselineskip}%
}

\let\lines\baselineskip

\begin{document}

\frontmatter

\maketitlepage

%\abstract{This is an example of {\tt infthesis} style.  The file {\tt
%skeleton.tex} generates this document and can be used to get a ``skeleton''
%for your thesis.  The abstract should summarise your report and fit in the
%space on the first page.
%%
%You may, of course, use any other software to write your report, as long as
%you follow the same style. That means: producing a title page as given here,
%and including a table of contents and bibliography.  }

\begin{publicationmeta}
  \section*{Acknowledgements}
  % TODO: Think of people to acknowledge I'd like to acknowledge Jesus for
  % swapping bottles whilst I was drinking them. It was really difficult to find
  % him at the bottom, so I had to keep drinking.

  \section*{Declaration}
  I declare that this thesis was composed by myself,
  that the work contained herein is my own
  except where explicitly stated otherwise in the text,
  and that this work has not been submitted for any other degree or
  professional qualification except as specified.\par
  ({\textit \thanklessauthor})
\end{publicationmeta}

\tableofcontents

%\pagenumbering{Arabic}

\mainmatter

\chapter{Introduction}\label{cha:introduction}

\section{Motivation}\label{sec:motivation}

In 1874 --- 4 years after his death --- John Papworth's \textit{Ordinary of British Armorials} was
published~\autocite{collins_1942}. In this work, he recorded approximately 50,000 entries of
descriptions of families' coats of arms, none annotated.

This honours project would make it possible to have these descriptions, or \textit{blazons} as they
are termed in heraldry (see~\ref{sec:heraldry}), drawn freely for people to view. This has
potential application for ancestry companies that build family trees for people. Given the blazon,
they would be able to construct the shields visually.

\section{Contributions}\label{sec:contributions}

In this honours project, my contributions included:

\begin{itemize}
  \item Drawing many of the charges used in the coats of arms
  \item Writing the base web server
  \item Writing the charge drawing algorithm
\end{itemize}

\chapter{Background}\label{cha:background}

\section{Heraldry}\label{sec:heraldry}

Many families, countries and organisations --- primarily in Europe --- have coats of arms. At the
centre of a coat of arms is a shield known as an \textit{escutcheon}. The language used to describe
how the escutcheon is to be drawn is known as a \textit{blazon}. Blazons have been used since the
Norman conquest and have been refined to a regular language in the process~\autocite{boutell_1864},
although, as John Brooke-Little said, ``many of the supposedly hard and fast rules laid down in
heraldic manuals [including those by heralds] are often ignored.''~\autocite{brooke_little_1985}
This flagrant disregard for the rules introduces difficulty in parsing the blazons as the language
loses some of its regularity.

Blazons have a few key attributes:
\begin{itemize}
  \item The \textit{field}, which is the background of the shield;
  \item \textit{Variations}, which describe how the field or charge are patterened. Variations can
    indicate patterns such as chequered or coloured lines (an example can be seen in
    Figure~\ref{fig:blason_albert});
  \item \textit{Ordinaries}, which are geometric shapes (an example of which can be seen in
    Figure~\ref{fig:scrope});
  \item \textit{Charges}, which are small emblems, such as fleur-de-lis and lions; and
  \item \textit{Tinctures}, which are the colours and patterns for charges, ordinaries and fields.
\end{itemize}

\begin{marginfigure}
  \centering
  \def\svgwidth{0.8\linewidth}
  \input{graphics/Blason_Albert.pdf_tex}
  \caption{The shield of the town of Albert, France. \textit{Barry of ten argent and
  gules}. Source:~\url{https://en.wikipedia.org/wiki/File:Blason_Albert.svg}}\label{fig:blason_albert}
\end{marginfigure}

The tinctures are derived from Norman French and are divided into 3 groups, typically known as
\textit{metals}, \textit{colours} and \textit{furs}. In British heraldry, the colours are derived
from Norman French and so the names can look archaic. In heraldry, blue is \textit{azure} and red
is \textit{gules} for instance. The metals are \textit{or} and \textit{argent} for gold and silver
respectively. Whilst the tinctures are named explicitly, the College of Arms does not define fixed
colours for the tinctures, leaving it to the artist to decide~\autocite{college_of_arms_faq}.

Blazons conventionally follow a form of starting with the \textit{tincture} or \textit{variation}
of the field. After the description of the field, \textit{ordinaries} (geometric shapes as seen in
Figure~\ref{fig:scrope}) and \textit{charges} (small emblems, such as fleur-de-lis and lions) are
named with their tinctures.

% Should probably take out the history lesson
A simple --- but notable --- blazon is that of the Scrope family. In the 14th century, the Baron
Scrope brought a case action against Sir Roberts Grosvenor when he noticed that they both had the
same coat of arms. Many witnesses gave evidence in the case, including Geoffrey
Chaucer~\autocite{scrope_grosvenor}. The case was ultimately decided in Scrope's favour. The Scrope
coat of arms has a blazon of \textit{Azure, a bend Or}; a depiction of this as drawn by the web app
written for this project can be seen in Figure~\ref{fig:scrope}.

\begin{marginfigure}
  \centering
  \def\svgwidth{0.8\linewidth}
  \input{graphics/scrope.pdf_tex}
  \caption{The Scrope escutcheon; \textit{Azure, a bend Or}.}\label{fig:scrope}
\end{marginfigure}

%\pagebreak

Whilst the Scrope arms are prominent in heraldry, they are simplistic and indicative of
medi\ae{}val arms. Coats of arms became more complex as they developed through the centuries, with
instances of quarterly shields, grand-quarterlies --- quarterlies within quarterlies --- and
\textit{differenced} arms. Examples of more complex shields can be seen in
Figure~\ref{fig:complex_shields}.

\begin{figure*}[h]
  \subfloat[Neville, 16th Earl of Warwick's coat of arms. An example of grand-quarterlies and
  differenced arms. Source:~\url{https://en.wikipedia.org/wiki/File:Neville_Warwick_Arms.svg}.]{%
    \def\svgwidth{0.3\linewidth} %
    \input{graphics/Neville_Warwick_Arms.pdf_tex}
  }
  \qquad
  \subfloat[Arms of Leveson-Gower, Earl Granville. An example of quarterlies with differenced
  arms. Source:~\url{https://en.wikipedia.org/wiki/File:LevesonGowerEarlGranvilleArms.png}.]{%
    \includegraphics[width=0.3\linewidth]{LevesonGowerEarlGranvilleArms.png}
  }
  \caption{Some examples of more complex coats of arms.}\label{fig:complex_shields}
\end{figure*}

For a time, it was considered bad form to repeat a \textit{tincture} in a blazon, and use a
reference to the tincture's previous use. The Heraldic Society gives an example as such:
``\textit{`Azure on a fess argent three billets azure'} [would have been written as] \textit{`Azure
on a fess argent three billets of the first'}''. The `\textit{of the first}' refers to the field's
tincture of azure. This blazon describes a blue shield, with a white bar horizontally across the
middle with 3 white rectangles arranged along the bar. The Heraldic Society advocates in favour of
repeating tinctures to reduce ambiguity~\autocite{blazon_in_coa}.

\backmatter

\printbibliography[heading=bibintoc]

\end{document}
